% Luận văn thạc sĩ KHMT
% Dựa trên Template KLTN cho SV trường ĐHKHTN
% Tác giả: Thầy NQMinh  (nqminh@fit.hcmus.edu.vn)  - 2016
%               BHThong (bhthong@fit.hcmus.edu.vn) - 2020
%               THQuan  (thquan@fit.hcmus.edu.vn)  - 2026
%
%
% Hướng dẫn: Điền thông tin vào khung CHÚ Ý bên dưới.
% Build: pdflatex > bibtex > pdflatex > pdflatex
%        Hoặc chạy: bash thesis-build.sh

\documentclass{hcmusthesis}

% ========================================================================================= %
% CHÚ Ý: Thông tin chung về KLTN - sinh viên điền vào đây để tự động update các trang khác  %
% Cú pháp lệnh thiết lập: \setThongTin{Tiếng Việt}{English}                                 %
% ========================================================================================= %
\setHoVaTen{Họ và Tên Học Viên}{Full Name of Student}               % Tên học viên: {Tiếng Việt}{English} (dùng ~ thay khoảng trắng)
\setMSSV{XXCXXXXX}                                                   % Mã số học viên
\setKhoaDaoTao{XX/20XX}                                              % Khóa đào tạo
\setTenKL{Tên~đề~tài~luận~văn~thạc~sĩ}{Thesis Title in English}      % Tên đề tài: {Tiếng Việt}{English} (dùng ~ thay khoảng trắng)
\setMaNganh{8.48.01.01}                                              % Mã số ngành (dùng chung cho cả tiếng Anh và tiếng Việt)
\setTenNganh{Tên ngành}{Speciality Name}                             % Tên ngành: {Tiếng Việt}{English}
\setTenGVHD{Học hàm. Họ và Tên GVHD}{Title. Full Name of Supervisor}% Người hướng dẫn: {Tiếng Việt}{English}
\setTenBM{Tên bộ môn}                                                % Tên bộ môn
\setDiaChiNgayThang{Tp. Hồ Chí Minh -- Năm 20XX}                     % Địa chỉ và ngày tháng bảo vệ

\hypersetup{
    pdfauthor={\tenSVTiengViet},
    pdftitle={\tenKLTiengViet},
    % pdfkeywords={TODO: Từ khóa},
    % pdfsubject={TODO: Tóm tắt}
}

\begin{document}

\pagenumbering{gobble}
\begin{titlepage}

\begin{center}
\MakeUppercase{đại học quốc gia tp. hồ chí minh}\\
\MakeUppercase{\bfseries trường đại học khoa học tự nhiên}\\[2cm]
% \textbf{\MakeUppercase{khoa công nghệ thông tin}}\\[2cm]


{ \large \bfseries \MakeUppercase{\tenSVTiengViet}\\[2cm] } 

%Tên đề tài Khóa luận tốt nghiệp/Đồ án tốt nghiệp

{ \Large \bfseries \MakeUppercase{\tenKLTiengViet} \\[3cm]} 


%Chọn trong các dòng sau
{\large\bfseries\MakeUppercase{luận văn thạc sĩ}}\\
%\large ĐỒ ÁN TỐT NGHIỆP CỬ NHÂN\\
%\large THỰC TẬP TỐT NGHIỆP CỬ NHÂN\\
%Đưa vào dòng này nếu thuộc chương trình Chất lượng cao, hoặc lớp Cử nhân tài năng
%\large CHƯƠNG TRÌNH CHÍNH QUY\\
%\large CHƯƠNG TRÌNH CHẤT LƯỢNG CAO\\
%\large CHƯƠNG TRÌNH CỬ NHÂN TÀI NĂNG\\[2cm]


\begin{tikzpicture}[remember picture, overlay]
  \draw[line width = 2pt] ($(current page.north west) + (3.25cm, -3.25cm)$) rectangle ($(current page.south east) + (-2.0cm, 3.0cm)$);
\end{tikzpicture}

\vfill
\diachiNgayThang

\end{center}

\pagebreak

\begin{center}
\MakeUppercase{đại học quốc gia tp. hồ chí minh}\\
\MakeUppercase{\bfseries trường đại học khoa học tự nhiên}\\[2cm]
% \textbf{\MakeUppercase{khoa công nghệ thông tin}}\\[2cm]


{\large \bfseries \MakeUppercase{\tenSVTiengViet}\\[2cm]} 


%Tên đề tài Khóa luận tốt nghiệp/Đồ án tốt nghiệp

{ \Large \bfseries \MakeUppercase{\tenKLTiengViet} \\[2cm]}  


%Chọn trong các dòng sau
% \large LUẬN VĂN THẠC SĨ KHOA HỌC MÁY TÍNH\\
%\large ĐỒ ÁN TỐT NGHIỆP CỬ NHÂN\\
%Đưa vào dòng này nếu thuộc chương trình Chất lượng cao, hoặc lớp Cử nhân tài năng
%\large CHƯƠNG TRÌNH CHÍNH QUY\\[2cm]
%\large CHƯƠNG TRÌNH CHẤT LƯỢNG CAO\\[2cm]
%\large CHƯƠNG TRÌNH CỬ NHÂN TÀI NĂNG\\[2cm]

\end{center}

Ngành: \tenNganhTiengViet

Mã số ngành: \maNganh \\[2cm]

\textbf{\MakeUppercase{người hướng dẫn khoa học}}
\begin{enumerate}[leftmargin=2.5cm]
    \item \tenGVHDTiengViet
\end{enumerate}

\begin{tikzpicture}[remember picture, overlay]
  \draw[line width = 2pt] ($(current page.north west) + (3.25cm, -3.25cm)$) rectangle ($(current page.south east) + (-2.0cm, 3.0cm)$);
\end{tikzpicture}

\vfill

\begin{center}
\diachiNgayThang
\end{center}
\end{titlepage}

\pagenumbering{roman} % Đánh số i, ii, iii, ...

\chapter*{Lời cam đoan}
\label{reassurances}

Tôi cam đoan luận văn thạc sĩ ngành \tenNganhTiengViet, với đề tài {\itshape\renewcommand{~}{ }\tenKhongXuongDong{\tenKLTiengViet}} là công trình khoa học do tôi thực hiện dưới sự hướng dẫn của \listSupervisors.

Những kết quả nghiên cứu của luận văn hoàn toàn trung thực và chính xác.\\[1em]

\hfill%
\begin{minipage}[r]{.5\textwidth}
    \centering
    Học viên cao học \\ [2.5cm]
    \tenSVTiengViet
\end{minipage}

\addcontentsline{toc}{chapter}{Lời cam đoan}

\chapter*{Lời cảm ơn}\label{thanks}

% ========================================================================================= %
% CHÚ Ý: Học viên viết lời cảm ơn ở đây.                                                    %
% - Lệnh \tenGVHDTiengViet (hoặc \listSupervisors) tự động cập nhật danh sách các CBHD       %
%   (đã định dạng tự động thêm dấu phẩy và từ "và" cho người cuối cùng).                    %
% ========================================================================================= %
% TODO: Điền lời cảm ơn vào đây.

Em xin cảm ơn thầy/cô hướng dẫn \listSupervisors\ vì những hướng dẫn và nhận xét quý báu cho luận văn này.

Em xin gửi lời cảm ơn đến quý thầy cô trong Khoa Công nghệ Thông tin, Trường Đại học Khoa học Tự nhiên -- ĐHQG HCM đã hỗ trợ em trong suốt quá trình học tập và nghiên cứu.

\addcontentsline{toc}{chapter}{Lời cảm ơn}

% Mục lục, danh sách hình, danh sách bảng
\tableofcontents
\addcontentsline{toc}{chapter}{Mục lục}
\listoffigures
\addcontentsline{toc}{chapter}{Danh mục các hình, biểu đồ}
\listoftables
\addcontentsline{toc}{chapter}{Danh mục các bảng số liệu}
\printnoidxglossary[type=\acronymtype, title=Danh~mục các ký~hiệu và từ~viết~tắt]

% Reset chapter number cho trang thông tin
\renewcommand{\thesection}{\arabic{section}}
\setcounter{section}{0}
\titleformat*{\section}{\large\bfseries}
\addtocontents{toc}{\protect\setcounter{tocdepth}{0}} 
\begin{trangthongtinluanvan}
    \begin{tomtatvi}
        % TODO: Học viên điền tóm tắt nội dung luận văn (Tiếng Việt) dưới đây.
        Luận văn giải quyết bài toán $\ldots$ Theo đó, luận văn đề xuất các hướng tiếp cận sau:

        \begin{enumerate}
            \item Hướng tiếp cận 1: $\ldots$
            \item Hướng tiếp cận 2: $\ldots$
        \end{enumerate}
    \end{tomtatvi}

    \begin{ketquamoivi}
        % TODO: Học viên điền các đóng góp/kết quả mới đạt được trong luận văn (Tiếng Việt) dưới đây.
        \begin{enumerate}
            \item Kết quả 1: $\ldots$
            \item Kết quả 2: $\ldots$
        \end{enumerate}
    \end{ketquamoivi}

    \begin{ungdungvi}
        % TODO: Học viên điền khả năng ứng dụng thực tiễn và hướng nghiên cứu tiếp theo (Tiếng Việt) dưới đây.
    \end{ungdungvi}
\end{trangthongtinluanvan}

\begin{trangthesisinformation}
    \begin{summaryen}
        % TODO: Student writes the thesis summary (English) below.
    \end{summaryen}

    \begin{noveltyen}
        % TODO: Student writes the novelty of the thesis (English) below.
    \end{noveltyen}

    \begin{applicationsen}
        % TODO: Student writes the practical applications and future research directions (English) below.
    \end{applicationsen}
\end{trangthesisinformation}


\clearpage

% Đánh số chapter.section từ trang này
\renewcommand{\thesection}{\thechapter.\arabic{section}}
\titleformat*{\section}{\Large\bfseries}
\pagenumbering{arabic} % Đánh số 1, 2, 3, ...
% Đếm đến subsection
\addtocontents{toc}{\protect\setcounter{tocdepth}{2}} 

% Các chương nội dung
\chapter{Giới thiệu}
\label{Chapter1}

% TODO: Giới thiệu tổng quan đề tài.

\section{Đặt vấn đề}

Đặt vấn đề và bối cảnh nghiên cứu \cite{adams2023}.

\section{Mục tiêu nghiên cứu}

Mục tiêu và phạm vi của luận văn.

\section{Đóng góp của luận văn}

\begin{enumerate}
    \item Đóng góp 1.
    \item Đóng góp 2.
\end{enumerate}

\section{Bố cục luận văn}

Luận văn được tổ chức như sau:

\begin{itemize}
    \item \textbf{Chương 1}: Giới thiệu.
    \item \textbf{Chương 2}: Các nghiên cứu liên quan.
    \item \textbf{Chương 3}: Phương pháp đề xuất.
    \item \textbf{Chương 4}: Thực nghiệm và đánh giá.
    \item \textbf{Chương 5}: Kết luận.
\end{itemize}


\chapter{Các nghiên cứu liên quan}
\label{Chapter2}

% TODO: Trình bày các nghiên cứu liên quan.

\section{Nghiên cứu liên quan 1}

Mô tả nghiên cứu liên quan thứ nhất \cite{sample2024}.

\section{Nghiên cứu liên quan 2}

Mô tả nghiên cứu liên quan thứ hai \cite{english_sample2025}.

\section{Tổng kết chương}

Tổng kết các nghiên cứu liên quan và định hướng của luận văn.


\chapter{Phương pháp đề xuất}
\label{Chapter3}

% TODO: Trình bày phương pháp đề xuất.

\section{Tổng quan phương pháp}

Tổng quan phương pháp đề xuất. Hình~\ref{fig:sample} minh họa kiến trúc tổng thể của hệ thống.

\begin{figure}[ht]
    \centering
    \includegraphics[width=0.7\textwidth]{sample-figure}
    \caption{Minh họa kiến trúc tổng thể của hệ thống.}
    \label{fig:sample}
\end{figure}

\section{Chi tiết phương pháp}

Chi tiết các bước trong phương pháp.

\section{Tổng kết chương}

Tổng kết phương pháp đề xuất.


\chapter{Thực nghiệm và đánh giá}
\label{Chapter4}

% TODO: Trình bày thực nghiệm và kết quả.

\section{Dữ liệu thực nghiệm}

Mô tả tập dữ liệu sử dụng trong thực nghiệm. Bảng~\ref{tab:dataset} trình bày thống kê dữ liệu.

\begin{table}[ht]
    \centering
    \caption{Thống kê tập dữ liệu.}
    \label{tab:dataset}
    \begin{tabular}{lccc}
        \toprule
        \textbf{Tập} & \textbf{Train} & \textbf{Dev} & \textbf{Test} \\
        \midrule
        Nhãn A       & 1{,}000         & 200           & 200           \\
        Nhãn B       & 800             & 150           & 150           \\
        Nhãn C       & 600             & 100           & 100           \\
        \midrule
        \textbf{Tổng} & \textbf{2{,}400} & \textbf{450} & \textbf{450} \\
        \bottomrule
    \end{tabular}
\end{table}

\section{Kết quả thực nghiệm}

Trình bày và phân tích kết quả thực nghiệm.

\section{Tổng kết chương}

Tổng kết kết quả thực nghiệm.


\chapter{Kết luận}
\label{Chapter5}

% TODO: Kết luận và hướng phát triển.

\section{Kết luận}

Luận văn đã trình bày $\ldots$

\section{Hướng phát triển}

Các hướng phát triển trong tương lai bao gồm $\ldots$


% Công trình của tác giả (nếu không có thì comment 02 dòng dưới)
\chapter*{Danh mục công trình của tác giả}\label{PublishedPapers}
\section*{Danh sách bài báo được chấp nhận trình bày ở các hội nghị}
\begin{enumerate}
% TODO: Thêm danh sách công trình, hoặc xóa chapter này nếu chưa có.
\item[CT-01] \hypertarget{pub:CT-01}{} \textbf{Họ Tên Tác Giả}, \textit{Tên bài báo}. Tên hội nghị, Năm.
\end{enumerate}

\addcontentsline{toc}{chapter}{Danh mục công trình của tác giả}

% In tài liệu tham khảo
\printbibheading[title={Tài liệu tham khảo}]
\addcontentsline{toc}{chapter}{Tài liệu tham khảo}

\printbibliography[heading=subbibliography, title={Tiếng Việt}, keyword=Viet, resetnumbers=true]

\DeclareNameAlias{sortname}{family-given}
\DeclareNameAlias{default}{family-given}

\printbibliography[heading=subbibliography, title={Tiếng Anh}, notkeyword=Viet, resetnumbers=1] 
% ===================================================================== %
% CHÚ Ý: phải gán lại resetnumbers=số tài liệu tham khảo tiếng Việt + 1 %
% ===================================================================== %

% Phần phụ lục
\appendix
\chapter{Phụ lục A}
\label{appendix:a}

% TODO: Điền nội dung phụ lục vào đây.

Đây là phụ lục ví dụ.


\end{document}
